\begin{definition} $\sat = \{ \varphi | \varphi \text{~--- выполнимая, то есть } \exists x: \varphi(x) = 1\}$
\end{definition}

\begin{definition} $3\sat = \{ \varphi | \varphi \text{~--- записана в 3-КНФ и выполнима} \}$
\end{definition}

\begin{note} Если $A \in \npc$ и $A \leqslant_p B$,  $B \in \np$, то $B \in \npc$.
\end{note}

\begin{theorem} (Кука-Левина) Задача $\sat$ является \npc.
\end{theorem}

\begin{proof}
Сначала покажем, что $\sat \in \np$. Сертификатом $s$ может служить выполняющий набор, а $V(\varphi,s) = \varphi(s)$ ~--- верификатором, который проверяет выполнимость формулы на наборе $s$.

Пусть $L \in \np$. Покажем, что $L \leqslant_p \sat$. Пусть $V$, работающий с сертификатами $s$ длины $q(n)$ ~--- верификатор из определения $L \in \np$. Будем так же считать, что $V$ задан МТ с одной, бесконечной вправо лентой и работающей не дольше $p(n)$ шагов. Идея состоит в записи булевой формулы, которая моделирует работу этой машины.

Напомним, что конфигурацией называется слово вида $aqb$, где $q \in Q$, а $a,b \in \Gamma^*$. Имеется в виду, что машина находится в состоянии $q$ и указывает на первый символ слова $b$, а левее записано слово $a$. Все остальные ячейки заполнены бланками. Вначале машина находится в конфигурации $q_0x\#s$, а в конце должна быть в состоянии $q_a$, если сертификат $s$ верный. Поскольку за 1 шаг машина сдвигается не больше, чем на одну ячейку вправо, максимальная длина любой конфигурации не больше $p(n)$. Введем служебные символы $\blacktriangleright$ и $\blacktriangleleft$, ограничивающие рабочую часть ленты. В новых обозначениях начальная конфигурация запишется как $\blacktriangleright q_0x\#s\#^{p(n)-n-q(n)-2}\blacktriangleleft$. 
Все возможные символы (элементы $\Gamma$, элементы $Q$, $\blacktriangleright$ и $\blacktriangleleft$) закодируем в двоичной записи, пусть на это потребуется $k$ битов. В таком случае для кодирования всех конфигурация нам потребуется $k(p(n)+2)(p(n) + 1)$ битов: $k$ битов на каждый символ, $(p(n)+2)$ символов в каждой конфигурации и $(p(n)+1)$ конфигураций за $p(n)$ шагов, включая начальную и конечную.

Теперь нужно построить формулу, гарантирующую корректность протокола. Она будет конъюнкцией трех формул, гарантирующих корректность начала, окончания и каждого шага. Для корректности начала нужно гарантировать, что первый символ равен $\blacktriangleright$, следующий ~--- $q_0$, следующие $n$ символов совпадают с битами $x$, далее идет бланк, далее $q(n)$ битов (т.е. неслужебных символов) далее бланки в количестве $p(n)-n-q(n)-2$ и, наконец, $\blacktriangleleft$. Каждое такое равенство выражается простой формулой, содержащей $k$ эквиваленций (в случае битов сертификата ~--- дизъюнкций двух наборов из $k$ эквиваленций), таким образом, общая длина конъюнкции всех формул будет порядка $p(n)$.

Для корректности окончания нужно построить формулу, проверяющую, что среди символов последней строки есть $q_a$. Это делается при помощи дизъюнкции эквиваленций. Наконец, нужно построить формулу, проверяющую корректность каждого шага. Ключевой идеей здесь является то, что вычисления на машине Тьюринга локальны, т.е. на каждом шаге изменяется небольшая часть ленты, и это изменение зависит от небольшой части ленты на предыдущем шаге. Рассмотрим таблицу, где в каждой строке записана конфигурация МТ на данном этапе. Тогда значение символа в каждой ячейке зависит от значений символов в $4$ ячейках
на предыдущем шаге: 
\resizebox{2cm}{!}{
\begin{tabular}{|c|c|c|c|}
    \hline
    ~ & ~ & ~ & ~ \\ \hline
    \multicolumn{1}{c|}{} & ~ & \multicolumn{1}{c}{} & \multicolumn{1}{c}{} \\ \cline{2-2}
\end{tabular}
} (зависимость от самой правой ячейки появляется, когда машина сдвигается налево: например, 
\resizebox{2cm}{!}{
\begin{tabular}{|c|c|c|c|}
    \hline
    \textbf{c} & \textbf{d} & \textbf{q} & \textbf{a} \\ \hline
    c & \textbf{r} & d & b \\ \hline
\end{tabular}
} для команды $qa \rightarrow rbL$). Вид этой зависимости полностью определяется программой машиной Тьюринга, а значит, может быть описан некоторой фиксированной пропозициональной формулой. Эту формулу нужно скопировать $p(n)(p(n) - 1)$ раз для каждого из возможных положений фигуры 
\resizebox{2cm}{!}{
\begin{tabular}{|c|c|c|c|}
    \hline
    ~ & ~ & ~ & ~ \\ \hline
    \multicolumn{1}{c|}{} & ~ & \multicolumn{1}{c}{} & \multicolumn{1}{c}{} \\ \cline{2-2}
\end{tabular}
} и взять конъюнкцию всех формул.

Таким образом, мы построили формулу, размер которой есть полином от исходных данных. Её выполнимость равносильна тому, что для некоторого сертификата $s$ машина на входе $x$ остановится в принимающем состоянии, т.е. $x \in L$. Такм образом, $L \leqslant_p \sat$ \end{proof}

Полезно будет так же следующее утверждение

\begin{statement} $\sat \leqslant_p 3\sat$
\end{statement}

\begin{proof}
Пусть $\varphi$ ~--- пропозициональная формула. Введём переменную для каждой её подформулы (включая исходные переменные и всю формулу). Каждая подформула, составленная из двух, например, задаёт утверждение вида $q = r \wedge s$. Его можно представить как 3-КНФ (в данном случае $(q \vee \neg r \vee \neg s) \wedge (\neg q \vee r \vee \neg s) \wedge (\neg q \vee \neg r \vee s) \wedge (\neg q \vee r \vee s)$). Взяв конъюнкцию всех таких формул, мы получим 3-КНФ, выполнимость которой эквивалентна выполнимости исходной формулы. Действительно, выполняющий набор исходной формулы задаст значения всех подформул, которые будут выполняющим набором полученной формулы. И наоборот, из выполняющего набора новой формулы можно выделить выполняющий набор исходной.

Осталось пояснить, почему этот алгоритм работает полиномиальное время. Ясно, что достаточно построить за полиномиальное время дерево синтаксического разбора формулы, дальнейшие операции занимают константное время для каждой его вершины. Построение дерева также несложно: путём подсчёта скобочного итога можно выделить из формулы две подформулы, из которых она получена, и повторить процедуру с ними рекурсивно. Это займёт линейное время для каждой подформулы, т.е. всего квадратичное время 
\end{proof}